\section{Introduction}\label{sec:intro}
False stories travel farther and faster than true ones on social platforms \citep{vosoughi2018spread}, and the cost of that asymmetry is borne by public health, elections and trust in institutions \citep{lazer2018science}. Automatic detection is one response, and on the standard benchmarks it appears close to solved: content-based classifiers routinely report accuracies above 95\% \citep{shu2017survey,zhou2020survey,kaliyar2021fakebert}. A detector, however, is deployed on data it was not trained on, and a growing body of work shows that these scores do not survive a change of corpus \citep{perezrosas2018automatic,janicka2019crossdomain,suprem2022generalizability,delloglio2026comparison} and that some of them reflect dataset artefacts rather than veracity \citep{asr2019bigdata,thibault2025guide,verma2026audit}.

What is still missing is an account of \emph{why} the transfer gap arises in a form that tells a practitioner what to do about it. A detector can fail on a new corpus in two different ways. It can \emph{order} the new items badly, scoring fabricated stories no higher than genuine ones; or it can order them well but place its decision threshold in the wrong place, because the new corpus has a different base rate or because its scores are scaled differently. The first failure needs a better representation; the second only needs a better operating point. Accuracy and F1, the metrics almost every cross-domain study reports, conflate the two.

This paper separates them. We assemble ten public corpora in two languages that cover four text units (articles, headlines, claims and tweets) and three constructs (veracity, rumour and clickbait), audit them for label direction and for within- and cross-corpus duplication, and evaluate six detectors in a train-on-one, test-on-all design with five seeds, bootstrap confidence intervals and permutation controls. We then decompose the macro-F1 that a detector loses on a new corpus into a ranking deficit, which no threshold can recover, and operating-point losses, which the right threshold would, and we evaluate label-free repairs aimed at each part. Our contributions are:

\begin{enumerate}
\item \textbf{An audited multi-corpus benchmark.} The two most used article corpora code their labels in opposite directions (ISOT: $1=$ fake; WELFake: $1=$ real). We show that simple signature probes detect such an inversion, which otherwise masquerades as a below-chance collapse of transfer, as it did in our own first pipeline. MinHash deduplication finds label-conflicting duplicates in GossipCop and source tweets shared between PHEME and Twitter15/16 (Section~\ref{sec:data}).
\item \textbf{A transfer gap decomposition.} Across 42 directed corpus pairs, \RankShareAll{} of the macro-F1 gap is a ranking deficit and \OpShareAll{} an operating-point loss (Section~\ref{sec:where}). Mean transfer AUC is \OffAUC{} against \InAUC{} in-domain; the gap is real but is not the collapse to chance that F1 alone suggests.
\item \textbf{A label-free operating-point repair, and why the standard one fails.} Expectation--maximisation prior adaptation, the textbook remedy for label shift, enlarges the operating-point loss (\RecEM{} recovered), because calibration does not survive the change of corpus. Transferring the ranking rule instead (flagging the source's base rate of target items) recovers \RecRank{} without any target label (Section~\ref{sec:op}).
\item \textbf{Evidence on what the detectors learn.} A model of 22 surface statistics and group identity alone recover much of the in-domain signal, 64 examples suffice on the article and tweet corpora, and the strongest cues are outlet signatures and time expressions; masking, feature deletion and topic balancing show how much of the ranking survives without them (Section~\ref{sec:shortcuts}). For the ranking deficit, training on several corpora at once raises held-out AUC from \LocoMeanSingle{} to \LocoPooled{}, whereas making the representation invariant to corpus identity (feature deletion, CORAL, nullspace projection) lowers it, because corpus identity and label are entangled in these benchmarks (Section~\ref{sec:repair-res}).
\item \textbf{A low-resource extension.} In Bengali, clickbait detection replicates the published baseline, survives held-out publishers and transfers across languages by translation, which suggests that whether a cue transfers depends on whether it \emph{is} the construct or merely correlates with it (Section~\ref{sec:bangla}).
\item \textbf{A reporting protocol and a full release} of code, seeds, per-cell results and predictions (Section~\ref{sec:discussion}).
\end{enumerate}
